\section{State Construction from Traces}
\label{sec:state_construction}

The formalism is useful only if traces can produce the fitted chain
reproducibly and with checks that expose when the abstraction fails. A
fitted Markov model is an estimate plus evidence: the trace
representation, order test, first-passage fit, and uncertainty analysis
must all support its use. Figure~\ref{fig:design} summarizes the
workflow; Algorithm~\ref{alg:trace} implements the construction pipeline
and Algorithm~\ref{alg:gof} the goodness-of-fit (GoF) protocol that SS7
(\S\ref{sec:ss7}) checks on controlled corpora.

%% [REDUCTION] The standalone pipeline figure is superseded by the study-design
%% figure (Fig.~\ref{fig:design}) in \S\ref{sec:intro}. Retained for the artifact.
\iffalse
\begin{figure}[!t]
\centering
\begin{tikzpicture}[
    >=Stealth, font=\scriptsize\sffamily, node distance=2pt,
    stage/.style={rectangle, rounded corners, draw=blue!65!black, fill=blue!6,
      line width=0.7pt, align=center, text width=1.55cm, minimum height=0.95cm, inner sep=2pt},
    gate/.style={diamond, aspect=1.7, draw=red!60!black, fill=red!7, line width=0.7pt,
      align=center, font=\scriptsize\sffamily, inner sep=1pt},
    result/.style={rectangle, rounded corners, draw=green!55!black, fill=green!7,
      line width=0.7pt, align=center, text width=1.7cm, minimum height=0.95cm, inner sep=2pt},
    ar/.style={->, draw=gray!75, line width=0.7pt}]
  \node[stage] (tr) {agent\\traces $\mathcal T$};
  \node[stage, right=10pt of tr] (fe) {featurize\\$\phi$ + cluster ($m$ states)};
  \node[stage, right=10pt of fe] (mle) {Laplace MLE\\$\hat Q,\hat R_\oplus,\hat R_\ominus$};
  \node[gate, below=12pt of mle] (gof) {AIC\\$\wedge$\,KS?};
  \node[result, left=14pt of gof] (use) {$\mathcal R(d)$, $R_\infty$,\\perturbation,\\pass$@k$/pass$^k$\\+ CIs (UQ)};
  \node[stage, right=12pt of gof, text width=1.6cm, draw=orange!70!black, fill=orange!8] (rej) {reject:\\re-featurize /\\richer model};
  \draw[ar] (tr) -- (fe);
  \draw[ar] (fe) -- (mle);
  \draw[ar] (mle) -- (gof);
  \draw[ar] (gof) -- node[above, font=\scriptsize] {pass} (use);
  \draw[ar] (gof) -- node[above, font=\scriptsize] {fail} (rej);
\end{tikzpicture}
\caption{\textsc{TraceToChain} workflow: traces are featurized and clustered
into $m$ transient states; Laplace-smoothed MLE fits the absorbing chain; the
composite AIC$\,\wedge\,$KS certificate gates whether the first-passage
quantities (with uncertainty) may be reported, or the corpus is referred back
for re-featurization or a richer model.}
\label{fig:pipeline}
\end{figure}
\fi


\subsection{Featurization and Discretization}
\label{sec:featurization}

The first step is to turn heterogeneous trace events into countable
states: reasoning steps, tool calls, and observations are featurized
and then clustered into labels such as ``plan'' or ``tool result.'' Let
$\mathcal{T} = \{\tau^{(i)}\}_{i=1}^{n}$ be the trace corpus. A run is
$\tau^{(i)} = (\sigma_1^{(i)}, \sigma_2^{(i)}, \ldots,
\sigma_{L_i}^{(i)}, y^{(i)})$, where each $\sigma_t^{(i)}$ is one
observed step and $y^{(i)} \in \{\oplus, \ominus\}$ is the terminal
outcome. A featurizer $\phi : \sigma \mapsto \mathbb{R}^p$ maps each
step to a vector representation, either:
\begin{itemize}
  \item \emph{rule-based} (tool type, retry flag, error code, used in
        our simulation study SS7 and MAST experiments), or
  \item \emph{learned} (a pretrained sentence encoder applied to the
        agent's chain-of-thought and interchangeable in
        Algorithm~\ref{alg:trace}).
\end{itemize}
We cluster the feature vectors $\{\phi(\sigma_t^{(i)})\}$ using
agglomerative Ward linkage~\cite{ward1963}. The number of clusters
$m \in [k_{\min}, k_{\max}]$ is selected by silhouette
score~\cite{rousseeuw1987}; we use $[3,6]$ on the controlled corpora of
SS9 and $[2,10]$ on the real corpora, with SS11 varying $k_{\max}$
explicitly. Each step is assigned to the selected
cluster, yielding
$s_1^{(i)}, \ldots, s_{L_i}^{(i)} \in \{1, \ldots, m\}$ followed by
the terminal outcome.

\subsection{Transition Estimation and Order Selection}
\label{sec:order_selection}

Estimating the fitted chain is then a counting problem with a direct
operational reading: a row of $\hat Q$ says where the agent goes next
if it has not terminated, and the matching entries of $\hat R_\oplus$
and $\hat R_\ominus$ say how often that state exits to success or fatal
failure. We fit these quantities by
Laplace-smoothed maximum-likelihood estimation (MLE) with
$\alpha=1$, so
$\hat{Q}_{ij} = (c_{ij} + \alpha) / (c_i + \alpha(m+2))$
where $c_{ij}$ is the count of transitions from state $i$ to state
$j$, $c_{i,\oplus}$ and $c_{i,\ominus}$ count terminal exits from
state $i$, and
$c_i = \sum_j c_{ij} + c_{i,\oplus}+c_{i,\ominus}$. The same
denominator normalizes the success and failure exits. Smoothing
prevents sparse rows from assigning zero probability to unobserved but
plausible exits.

We also test whether first-order memory is sufficient. The comparison
uses a first-vs-second-order Markov Akaike information criterion (AIC)
test~\cite{akaike1974}:
$\Delta_{\rm AIC} = \mathrm{AIC}_1 - \mathrm{AIC}_2
                  = -2(\ell_1 - \ell_2) + 2(k_1 - k_2),$
where $\ell_o$ and $k_o$ are the log-likelihood and free-parameter
count under order $o$. The first-order model is preferred iff
$\Delta_{\rm AIC} < 0$. Simulation study SS7 (\S\ref{sec:ss7}) shows
that this test has high power against second-order ground truth.

\paragraph{Time-homogeneity as a testable assumption.}
When we pool transition counts, we assume a \emph{time-homogeneous}
kernel $\Pr[s_{t+1}=j \mid s_t=i]$ that does not vary with $t$. The
composite AIC$\,\wedge\,$KS protocol (Algorithm~\ref{alg:gof}) tests
whether that simplification is defensible. AIC checks whether recent
history still matters after conditioning on the current state, while
the Kolmogorov--Smirnov (KS) first-passage test can reveal drift in
exit hazards. If either test fails, the corpus should be segmented
before counts are pooled.

\paragraph{Variable horizons, early termination, and censoring.}
Runs can have different lengths and may end in success ($\oplus$),
failure ($\ominus$), or right-censoring. Algorithm~\ref{alg:trace}
counts each observed transient transition once. Completed runs add one
absorber count, while censored runs contribute only the transitions
observed before censoring. This keeps $\hat Q$ tied to observed
movement among transient states, but it can make the exit estimates
$\hat R_\oplus$ and $\hat R_\ominus$ conservative when censoring is
substantial. In \S\ref{sec:mast} we observe censoring below $10\%$, and
SS12 (\S\ref{sec:robustness}) probes far heavier regimes.

\begin{algorithm}[t]
\small
\DontPrintSemicolon
\KwIn{Traces $\mathcal{T}$, featurizer $\phi$, cluster range
      $[k_{\min}, k_{\max}]$, Laplace $\alpha$.}
\KwOut{$\hat{Q}, \hat{R}_\oplus, \hat{R}_\ominus$, labels $\pi$,
       AIC order verdict.}
$X \leftarrow [\phi(\sigma_t^{(i)})]_{i,t}$\tcp*{featurize}
\For{$k = k_{\min}$ \KwTo $k_{\max}$}{
  $L_k \leftarrow$ AgglomerativeWard$(X, k)$;\quad
  $\mathrm{sil}_k \leftarrow $ silhouette$(X, L_k)$\;
}
$m^\star \leftarrow \arg\max_k \mathrm{sil}_k$;\quad
$\pi \leftarrow L_{m^\star}$\;
$c_{s_t,s_{t+1}}\mathrel{+}=1$ for each transient transition;
$c_{s_{L_i},y^{(i)}}\mathrel{+}=1$ for each terminal exit\;
$\hat{Q}_{ij} \leftarrow (c_{ij} + \alpha) / (c_i + \alpha(m^\star + 2))$;
same normalization for $\hat{R}_{\oplus,i},\hat{R}_{\ominus,i}$\;
Compute $\ell_1, \ell_2, k_1, k_2$;\quad
$\Delta_{\rm AIC} \leftarrow -2(\ell_1 - \ell_2) + 2(k_1 - k_2)$\;
\KwRet{$(\hat{Q}, \hat{R}_\oplus, \hat{R}_\ominus, \pi)$;
       first-order iff $\Delta_{\rm AIC} < 0$}\;
\caption{\textsc{TraceToChain}: construct and audit an absorbing DTMC
         from agent execution traces.}
\label{alg:trace}
\end{algorithm}

\subsection{Goodness-of-Fit Protocol}
\label{sec:gof}

Fitting helps only if the abstraction matches the traces. A rejection
is an engineering result, not a nuisance: it says the current state
labels should not be used for reliability queries without segmentation,
new features, or a richer model. We use two checks on $\hat M$ that
address different failure modes:

\emph{Layer 1 (KS on first-passage).} We compare the model
success first-passage-time (FPT) distribution
$F_{\tau_\oplus \mid \tau_\oplus < \tau_\ominus}^{\hat M}$, conditional
on eventual success, with the empirical distribution of success FPTs
among successful traces, and retain the null iff $p_{\rm KS} > 0.05$.
The reference distribution is drawn from $\hat M$ as $N$ seeded
absorption times ($N=8{,}000$ in SS9, $20{,}000$ on the real corpora),
so this is a genuine \emph{two-sample} KS test~\cite{smirnov1939}, not a
one-sample test against a closed-form CDF: stochastic, but
seed-reproducible.

\emph{Layer 2 (AIC).} The AIC layer asks whether the first-order
chain is adequate compared with a second-order alternative. We reject
the first-order model if $\Delta_{\rm AIC} > 0$.

\emph{Composite rule.} The chain is accepted iff
\emph{both} $p_{\rm KS} > 0.05$ \emph{and} $\Delta_{\rm AIC} < 0$.
This rule catches second-order dependence that the FPT-marginal KS test
can miss (see \S\ref{sec:ss7}). If either layer fails, later
reliability quantities should be treated as unsupported by the current
trace abstraction. Given $\phi$, the cluster range, $\alpha$, and a
lowest-$k$ silhouette tie-breaking rule, Algorithm~\ref{alg:trace} is
deterministic in the trace corpus, and Algorithm~\ref{alg:gof} is
deterministic given its reference-sample seed; all seeds are fixed in
the artifact.

\begin{algorithm}[t]
\small
\DontPrintSemicolon
\KwIn{Traces $\mathcal{T}$, fitted chain $\hat{M}$, threshold
      $\alpha_{\rm KS}=0.05$.}
\KwOut{\textsc{Accept} / \textsc{Reject}, $(p_{\rm KS},\Delta_{\rm AIC})$.}
Draw $N$ seeded absorption times from $\hat M$ $\to$ reference
$\tilde F_{\tau_\oplus \mid \tau_\oplus < \tau_\ominus}^{\hat M}$\;
Collect empirical FPTs from $\mathcal{T}_\oplus$ (success traces)\;
$p_{\rm KS} \leftarrow$ two-sample KS($\hat F$, $\tilde F^{\hat M}$)\;
$\Delta_{\rm AIC} \leftarrow$ (Algorithm~\ref{alg:trace}, line~8)\;
\eIf{$p_{\rm KS} > \alpha_{\rm KS}$ \textbf{and} $\Delta_{\rm AIC} < 0$}{
  \KwRet{\textsc{Accept}, $(p_{\rm KS}, \Delta_{\rm AIC})$}\;
}{
  \KwRet{\textsc{Reject}, $(p_{\rm KS}, \Delta_{\rm AIC})$}\;
}
\caption{Composite KS/AIC goodness-of-fit test for the agent-DTMC
         assumption.}
\label{alg:gof}
\end{algorithm}

\subsection{Uncertainty Quantification on \texorpdfstring{$\hat Q$, $\hat R_\oplus$, $\hat R_\ominus$}{Q, R+, R-}}
\label{sec:uq}

Algorithm~\ref{alg:trace} returns point estimates, but a small corpus
may estimate a row of $\hat Q$ or an exit probability poorly, and that
uncertainty should be visible before an operator interprets $R(d)$,
$R_\infty$, or a perturbation result. We therefore attach intervals to
$\hat Q$, $\hat R_\oplus$, and $\hat R_\ominus$ and propagate them to
the quantities of \S\ref{sec:core}--\S\ref{sec:ext}.

\paragraph{(i) Closed-form Dirichlet posterior.}
The Laplace-smoothed MLE is the posterior mean under a symmetric
Dirichlet$(\alpha)$ prior over $\{s_1,\ldots,s_m,\oplus,\ominus\}$. The
row-$i$ posterior is $\mathrm{Dir}(c_{i,\cdot} + \alpha \mathbf{1})$
and the marginal for any single entry $X_{i,j}$ (including absorber
columns) is $\mathrm{Beta}(c_{i,j}+\alpha,\; c_i + \alpha(m+2) -
c_{i,j} - \alpha)$, where $c_i = \sum_{j'} c_{i,j'}$ includes both
transient and absorber outcomes. Equal-tailed $(1-\alpha_{\rm CI})$
credible intervals (CIs) are the corresponding Beta quantiles; we use
$\alpha_{\rm CI}=0.05$ throughout.

\paragraph{(ii) Non-parametric trace-level bootstrap.}
To quantify trace-sampling variability we resample traces with
replacement $B=200$ times, assign each step to its nearest target
centroid to keep labels aligned, and recompute the MLE. The implemented full
re-clustering path (\texttt{fast=False}) gives comparable but wider
intervals.

\paragraph{Empirical ranges.}
Table~\ref{tab:mast-states} reports representative MAST-derived state
taxonomies from SS8; each cluster is named by its dominant rule-based
feature (\S\ref{sec:featurization}) with that feature's within-cluster
share in parentheses, and the posterior intervals show which
success exits are well estimated or sparse, exposing state-level
variation hidden by aggregate pass/fail rates. Across the seven frameworks the two
uncertainty routes agree closely on entries of $\hat Q$: median
posterior CI widths span $0.108$--$0.116$ and median bootstrap widths
$0.084$--$0.103$, while maxima reaching $0.34$ mark the sparse rows. All downstream $\mathcal{R}(d)$, mean time to
absorption ($\mathrm{MTTA}=e_{s_0}^{\top}N\mathbf{1}$), and pass$^k$
quantities inherit these CIs
by propagation (Lipschitz-continuous in $(\hat Q,\hat R_\oplus)$ by
Proposition~\ref{thm:perturb}).

% Auto-generated by SS8_uncertainty_quantification.py
\begin{table}[t]
\centering
\caption{MAST-derived state taxonomy with 95\% posterior success-exit intervals.}
\label{tab:mast-states}
\scriptsize
\setlength{\tabcolsep}{2pt}
\begin{tabular}{l l c c c}
\toprule
\textbf{Framework} & \textbf{Cluster} & \textbf{$\hat{R}_{\oplus,i}$} & \textbf{95\% post.\ CI} & \textbf{$\hat{R}_{\ominus,i}$} \\
\midrule
react~\cite{yao2023react} & tool\_call (100\%) & 0.064 & [0.039,0.094] & 0.067 \\
\tstriped
 & plan (100\%) & 0.076 & [0.048,0.109] & 0.072 \\
 & retry (100\%) & 0.034 & [0.004,0.092] & 0.085 \\
\tstriped
 & reflect (100\%) & 0.081 & [0.038,0.138] & 0.081 \\
 & error\_parse (100\%) & 0.082 & [0.028,0.162] & 0.082 \\
\tstriped
 & wait (100\%) & 0.057 & [0.007,0.153] & 0.057 \\
\midrule
reflexion~\cite{shinn2023reflexion} & plan (100\%) & 0.121 & [0.074,0.178] & 0.027 \\
\tstriped
 & error\_parse (100\%) & 0.083 & [0.041,0.138] & 0.033 \\
 & retry (100\%) & 0.054 & [0.020,0.103] & 0.045 \\
\tstriped
 & reflect (100\%) & 0.104 & [0.056,0.166] & 0.043 \\
 & tool\_call (100\%) & 0.097 & [0.066,0.133] & 0.064 \\
\tstriped
 & wait (100\%) & 0.128 & [0.049,0.236] & 0.064 \\
\bottomrule
\end{tabular}
\end{table}

%% [REDUCTION] tab:uq-summary is a 7x7 grid whose every range is quoted
%% inline below; retained in the artifact.
%% % Auto-generated by SS8_uncertainty_quantification.py
\begin{table}[t]
\centering
\caption{Per-framework 95\% uncertainty widths for entries of $\hat Q$.}
\label{tab:uq-summary}
\scriptsize
\setlength{\tabcolsep}{2pt}
\begin{tabular}{l c c c c c c}
\toprule
 & \multicolumn{3}{c}{\textbf{Posterior CI width}} & \multicolumn{3}{c}{\textbf{Bootstrap CI width}} \\
\cmidrule(lr){2-4}\cmidrule(lr){5-7}
\textbf{Framework} & \textbf{min} & \textbf{median} & \textbf{max} & \textbf{min} & \textbf{median} & \textbf{max} \\
\midrule
react~\cite{yao2023react} & 0.037 & 0.110 & 0.313 & 0.017 & 0.097 & 0.287 \\
\tstriped
reflexion~\cite{shinn2023reflexion} & 0.048 & 0.109 & 0.238 & 0.040 & 0.098 & 0.253 \\
cot\_agent~\cite{wang2024agentsurvey} & 0.029 & 0.111 & 0.337 & 0.008 & 0.084 & 0.288 \\
\tstriped
toolformer~\cite{schick2023toolformer} & 0.041 & 0.108 & 0.241 & 0.034 & 0.101 & 0.208 \\
babyagi~\cite{wang2024agentsurvey} & 0.040 & 0.113 & 0.286 & 0.035 & 0.100 & 0.291 \\
\tstriped
autogpt~\cite{wang2024agentsurvey} & 0.044 & 0.112 & 0.283 & 0.012 & 0.100 & 0.268 \\
agentbench~\cite{agentbench2024} & 0.034 & 0.116 & 0.265 & 0.032 & 0.103 & 0.250 \\
\bottomrule
\end{tabular}
\end{table}


The output is the audited model used for reliability analysis: a
fitted chain plus uncertainty and fit diagnostics. If either layer
rejects, the corpus should be segmented, re-featurized, or treated as
outside the absorbing-DTMC approximation; when both pass, the fitted
chain can be used for reliability queries with the reported intervals.

%% ================================================================
